\changecolours{ScoutGreen}
\section[Cache-Aside e TTL]{Padrão Cache-Aside e Estratégias de Expiração (TTL)}

% ------------------------------------------------------------------------------
\twocolframe[title={O Padrão Cache-Aside (Lazy Loading)},leftcol=7.1cm,rightcol=6.9cm]{%
  \alert{Mecânica de Funcionamento}\par
  \itemseps[topsepi=0.2mm,itemsepi=0.4mm,labelsep=1.5mm,leftmargini=2.5mm]
  \begin{itemize}
    \item \textbf{A Aplicação Coordena:} O back-end busca na memória rápida (Redis) e na base persistente (Disco/SQL).
    \item \textbf{Fluxo de Leitura:}
      \begin{enumerate}
        \item \textbf{Cache Hit:} Se a chave existir, retorna da RAM ($\approx 0.4\text{ ms}$).
        \item \textbf{Cache Miss:} Se não existir, busca na base lenta ($\approx 120\text{ ms}$), salva no Redis com \texttt{SETEX} e retorna.
      \end{enumerate}
    \item \textbf{Vantagem:} Apenas dados ativamente consultados consomem memória RAM.
  \end{itemize}
}{%
  \alert{Por que o Cache DEVE Ter TTL?}\par
  \itemseps[topsepi=0.2mm,itemsepi=0.4mm,labelsep=1.5mm,leftmargini=2.5mm]
  \begin{itemize}
    \item \textbf{Finitude de RAM:} Evita que o Redis atinja \texttt{maxmemory} e entre em evicção forçada.
    \item \textbf{Autocura:} Dados defasados (\textit{Stale Reads}) são eliminados após o tempo de vida.
    \item \textbf{Invalidação Pró-Ativa:} Ao atualizar o catálogo, deleta-se a chave (\texttt{DEL cache:catalogo}).
  \end{itemize}
}

% ------------------------------------------------------------------------------
\begin{frame}[fragile]{Implementação do Cache-Aside em Python no FastCommerce}
  \vspace{0.1cm}
  \begin{columns}[t]
    \begin{column}{0.48\textwidth}
      \footnotesize
      \alert{Lógica de Negócio na API}\par
      \vspace{0.15cm}
      \begin{itemize}\setlength{\itemsep}{1.2mm}\setlength{\parskip}{0pt}
        \item Utiliza \texttt{r.get(chave)} para consulta em memória.
        \item Grava no Redis com \texttt{r.setex(chave, 45, dados)} aplicando TTL de 45 segundos.
        \item Mensura a latência em tempo real para o painel de telemetria da interface web.
        \item Speedup comprovado: de $\approx 120\text{ ms}$ no Miss para $\approx 0.4\text{ ms}$ no Hit ($> 250\times$ mais rápido!).
      \end{itemize}
    \end{column}

    \begin{column}{0.48\textwidth}
      \alert{Código Python no FastCommerce}\par
      \vspace{0.1cm}
\begin{lstlisting}[language=Python,basicstyle=\fontsize{4.6pt}{5.6pt}\selectfont\ttfamily]
@app.route("/api/produtos")
def get_produtos():
    inicio = time.time()
    chave = "cache:catalogo_produtos"

    # 1. Verifica Cache no Redis (RAM)
    if r.exists(chave):
        produtos = json.loads(r.get(chave))
        origem = "CACHE_HIT (RAM Redis)"
    else:
        # 2. Cache Miss: Busca lenta simulada
        time.sleep(0.12)
        produtos = CATALOGO_DISCO
        # 3. Grava no Redis com TTL de 45s
        r.setex(chave, 45, json.dumps(produtos))
        origem = "CACHE_MISS (Disco/DB)"

    latencia = (time.time() - inicio) * 1000
    return jsonify({
        "origem": origem, 
        "latencia_ms": round(latencia, 2),
        "produtos": produtos
    })
\end{lstlisting}
    \end{column}
  \end{columns}
\end{frame}
